\section{Modelo de Métricas y Fórmulas Matemáticas}
\label{sec:metrics}

Las métricas mostradas en el Dashboard de Detection Attacks se derivan de la consolidación de los registros históricos en la tabla `validation\_logs`. Esta sección detalla la construcción conceptual y matemática de dichos indicadores.

\subsection{Criterio de Clasificación de Fraude}
El sistema asigna un veredicto binario sobre la veracidad del documento basándose en la puntuación de confianza $S$ retornada por el microservicio de AWS Rekognition. El modelo está calibrado para que $S \in [0, 1]$, donde un valor cercano a $0$ representa un documento real y un valor cercano a $1$ indica una sospecha o intento de fraude.

Definimos el veredicto de fraude como una función indicadora $I_{\text{fraude}}$ con un umbral de decisión $\tau = 0.5$:

\begin{equation}
I_{\text{fraude}}(S) = \begin{cases} 
1 & \text{si } S \ge 0.5 \quad \text{(FRAUDE / Sospecha)} \\
0 & \text{si } S < 0.5 \quad \text{(REAL / Aprobado)}
\end{cases}
\end{equation}

Este umbral coincide con la línea de decisión visible en los histogramas de distribución del dashboard.

\subsection{Indicadores Clave de Rendimiento (KPIs)}
Sea $N_{\text{total}}$ el número total de transacciones registradas en el webhook, y sea $M$ el subconjunto de transacciones correspondientes al análisis de imágenes frontales ($N_{\text{front}}$) o traseras ($N_{\text{back}}$):

\begin{equation}
M = \{x \in \text{logs} \mid \text{node}(x) \in \{\text{'process\_front'}, \text{'process\_back'}\} \}
\end{equation}

Definimos los indicadores de volumen del siguiente modo:
\begin{align}
\text{Total Validations} &= N_{\text{total}} \\
\text{Predicted Real} &= \sum_{i \in M} (1 - I_{\text{fraude}}(S_i)) \\
\text{Predicted Fraud} &= \sum_{i \in M} I_{\text{fraude}}(S_i)
\end{align}

Las tasas porcentuales de clasificación se formulan como:
\begin{equation}
\text{Percentage of Real (\%)} = \left( \frac{\text{Predicted Real}}{|M|} \right) \times 100
\end{equation}

\begin{equation}
\text{Percentage of Fraud (\%)} = \left( \frac{\text{Predicted Fraud}}{|M|} \right) \times 100
\end{equation}

\subsection{Histogramas de Distribución de Scores}
Para analizar la distribución probabilística del clasificador, se agrupan los scores continuos de cada transacción en $K=10$ bins de ancho constante $W = 0.1$. El conteo de registros $C_k$ en el bin $k$ (donde $k \in \{1, 2, \dots, 10\}$) se define como:

\begin{equation}
C_k = |\{i \in M \mid S_i \in [ (k-1) \cdot 0.1, \ k \cdot 0.1 ) \}|
\end{equation}
*(Para el último bin $k=10$, el intervalo es cerrado por la derecha $[0.9, 1.0]$).*

\subsection{Promedios de Latencia y Tiempos de Respuesta}
Sea $L_{\text{AWS}, i}$ la latencia en milisegundos de la transacción $i$ en AWS Lambda, y sea $L_{\text{Gemini}, i}$ la latencia del servicio de humanización de Gemini. Los promedios aritméticos generales mostrados en el dashboard se calculan como:

\begin{equation}
\bar{L}_{\text{AWS}} = \frac{1}{|M|} \sum_{i \in M} L_{\text{AWS}, i}
\end{equation}

\begin{equation}
\bar{L}_{\text{Gemini}} = \frac{1}{|M|} \sum_{i \in M} L_{\text{Gemini}, i}
\end{equation}
